% Suivi d'une réaction qui produit un gaz, en enceinte fermée (L'essentiel et exercice 10, chapitre 2)
\begin{tikzpicture}[x=1cm,y=1cm]
\pic at (0,0) {erlen={hauteur=2.4,niveau=0.8}};
\pic at (0.05,0.12) {bulles};
\filldraw[fill=otInk!65,draw=otInk!80,line width=0.4pt] (-0.3,2.4) -- (-0.26,2.85) -- (0.26,2.85) -- (0.3,2.4) -- cycle;
\draw[otInk!70,line width=1.1pt] (0.1,2.85) -- (0.1,3.35) -- (1.55,3.35);
\filldraw[pcAppareil,rounded corners=0.8mm] (1.55,3.05) rectangle (2.5,3.65);
\node[pcLegende] at (2.03,3.35) {capteur};
\draw[pcFil] (2.5,3.35) to[out=0,in=180] (3.25,2.4);
\pic (af) at (4.05,1.85) {afficheur={largeur=1.6,hauteur=1.1,texte={1352}}};
\node[pcLegende,anchor=north] at (4.05,1.8) {pression en hPa};
\draw[pcRepere] (-0.55,0.45) -- (-1.25,0.85) node[pcLegende,anchor=east]{réactifs\\en solution};
\draw[pcRepere] (-0.3,1.85) -- (-1.05,2.35) node[pcLegende,anchor=east]{gaz : volume $V$\\à température $T$};
\end{tikzpicture}
\par\smallskip\legende{À volume et température constants, l'augmentation de pression $\Delta p$ mesure la quantité de gaz formé : $n=\dfrac{\Delta p\,V}{RT}$.}
